\chapter{RSiena quirks reproduced on purpose}\label{app:quirks}

\nwsaom{} is built to give the same numbers as RSiena. Where RSiena's own
behaviour is odd, for example where it counts a tie twice or reads a covariate
at another node's index, the quirk is copied on purpose, because a different but
more sensible definition would give statistics and estimates that no longer
match \texttt{siena07}. This appendix lists the cases that matter to a user who
compares output of the two programs or reads the effect definitions closely.
Each was found by comparing the observed statistics (the targets) with those of
\texttt{siena07}, and each is validated that way.

\section*{Conventions that change the numbers}

\begin{itemize}
\item \textbf{Centring.} Constant actor covariates are centred by their mean
  (RSiena's \texttt{coCovar}), changing covariates by their overall mean over
  all waves. The degree-weighted and distance-two covariate families depend on
  this: with the raw covariate their statistics are off by the mean.
\item \textbf{Similarity.} The similarity covariate is
  $(x-\bar x)/\text{range}$ with the similarity mean as the centring constant.
  In the similarity effects of a behavior, equal alters join the \emph{higher}
  popularity count.
\item \textbf{Behavior variance.} The variance used for behavior effects is the
  pooled population variance of waves $1,\dots,W-1$, and the behavior mean is
  taken over all waves.
\item \textbf{Start network.} Statistics of behavior effects that involve the
  network use the network at the start of the period.
\item \textbf{Scaling of \texttt{gwdspfb}.} The decay used by the effect is
  $0.01$ times RSiena's parameter, so RSiena's 69 corresponds to a decay
  of 0.69.
\end{itemize}

\section*{Single effects}

\begin{itemize}
\item \textbf{\opt{reciact} (root version).} When an actor with no mutual tie
  withdraws a tie, RSiena evaluates the square-root table at $-1$, which reads
  the memory just before the table and yields about 0. \nwsaom{} takes the value
  to be 0.
\item \textbf{\texttt{allDifCovNetNetIn} and \texttt{contrastCovNetNetIn}} are
  the same function in RSiena and give the same statistic.
\item \textbf{Covariate stars.} \texttt{DifferentCovariateOutStar} tests the tie
  from the alter to $k$, whereas \texttt{SameCovariateOutStar} tests the tie from
  $k$ to the alter. Both are reproduced as they are.
\item \textbf{Average versions of the same-covariate effects} divide by the number
  of actors with the ego's covariate value, not by the number of alters.
\item \textbf{\opt{avdeg} in a two-mode network} divides by the number of
  \emph{actors}, not by the number of nodes.
\end{itemize}

\section*{Two-mode networks}

\begin{itemize}
\item \texttt{crossXOutAct} reads the actor covariate at the index of the
  \emph{event}, as RSiena does.
\item \texttt{avGroupEgoX} averages over the first $m$ actors and divides by
  $m$.
\item \texttt{homXOutAct2} has the default parameter 3 in a bipartite network.
\item \texttt{outInAss}/\texttt{outOutAss} and their relatives use the square
  roots that RSiena uses as the two-mode default (parameter 2).
\item In \texttt{nwsaom multinet}, \texttt{divInAltIntn}, \texttt{JoutMix} and
  \texttt{gwespFBMix} of a two-mode dependent network, and
  \texttt{outOutActIntn}/\texttt{outOutAvIntn} of a one-mode network that depends on a
  two-mode one, index the other network by the node number of the \emph{event}. \nwsaom{}
  copies this and refuses these effects when the events outnumber the actors.
\item \texttt{netbeh(totgroupegox)}, \texttt{netbeh(avgroupegox)} and
  \texttt{netbeh(crossxoutact)} in a two-mode network use the behavior of the actors for
  the events as well (the events mirror the actors' behavior), which needs at least as many actors
  as events.
\end{itemize}

\section*{Other effects}

\begin{itemize}
\item \textbf{\texttt{toU}.} In RSiena 1.6.12 and later \texttt{toU} counts, for the ego $i$
  and the alter $j$, the actors $q$ with a tie $i\to q$ in the \emph{other} network and a
  tie $q\to j$ in the \emph{own} network, each weighted by the dyadic covariate value $D(i,q)$.
  \texttt{sharedToU} and \texttt{avAltU} are read in the same way.
\item \textbf{\opt{higher} in a non-directed network} equals the density statistic.
  The effect is not identified; RSiena returns a standard error of NA and \nwsaom{} stops with
  an error message.
\item \textbf{The shape effect \texttt{constant}} of a behavior is confounded with the
  behavior rate (RSiena gives a standard error of NA), so a fit that includes it stops.
\item \textbf{The five double-degree behavior effects} (\texttt{FFDeg}, \texttt{BBDeg},
  \texttt{FBDeg}, \texttt{FRDeg}, \texttt{BRDeg}) together are not identified, in RSiena as well.
\item \textbf{\texttt{netbeh(indties)}} differs from \texttt{siena07} and is refused in
  \texttt{nwsaom multinet}.
\item \textbf{Internal parameters.} In the current RSiena \texttt{outPop.c}, \texttt{outPopSqrt},
  \texttt{inAct.c}, \texttt{inActSqrt}, \texttt{outPopMore} and relatives have the internal
  parameter 1 by default; \nwsaom{} follows.
\end{itemize}

\section*{Where \nwsaom{} deliberately differs}

\begin{itemize}
\item \textbf{\texttt{outPopSqrt} in older RSiena versions.} In RSiena 1.6.1 to 1.6.12 the
  default parameter~0 of \texttt{outPopSqrt} made the statistic use unrooted start-of-period
  outdegrees while the ministep still used the square root, and models with it did not converge.
  RSiena 1.6.13 changed the default to parameter~1. \nwsaom{} implements the documented effect, which
  equals the current RSiena default (the benchmark in Chapter~\ref{ch:benchmark} sets the parameter
  to 1 for the older version it used).
\item \textbf{Refusal instead of noise.} Some combinations are not reproduced
  because RSiena itself fails or returns noise: \texttt{sameXOutAct} and
  \texttt{diffXOutAct} in a two-mode network abort \texttt{siena07}; \texttt{inPop\_ego}
  in a two-mode network throws an exception in RSiena and \texttt{homXOutAct2} with a behavior is not implemented there;
  the non-directed \texttt{higher} effect is collinear with density. \nwsaom{}
  refuses these with an error message instead (Section~\ref{sec:cannoteval}).
\end{itemize}

